\chapter{Metodologia}
\label{cap:metodo}

A pesquisa compreende desenvolvimento, verificação funcional, calibração dos instrumentos e comparação experimental. A verificação do software identifica se cada mecanismo funciona nas condições testadas; a comparação exige o conjunto de repetições e controles especificado neste capítulo. Os registros da primeira etapa não serão utilizados como substitutos da coleta definitiva.

\section{Tipo de pesquisa}

\subsection{Natureza da pesquisa}
A pesquisa é aplicada porque constrói um protótipo para apoiar decisões de engenharia de software. O procedimento principal é experimental: duas configurações de comunicação serão submetidas a cargas e falhas controladas. Desempenho e recuperação dependem de condições de execução e, portanto, não podem ser estabelecidos apenas pela descrição das tecnologias. A definição prévia de variáveis e critérios de análise segue as recomendações de \citeonline{kitchenham2002}.

\subsection{Abordagem metodológica}
A abordagem é quantitativa. Latência, throughput, erros, recursos, backlog e recuperação serão expressos por medidas observáveis. A interpretação relacionará os números aos mecanismos executados e às limitações do ambiente. Cada execução completa constitui uma repetição; os eventos de uma mesma execução não serão tratados como repetições experimentais independentes.

A relação com os objetivos específicos organiza o procedimento. Contrato e regras comuns estabelecem equivalência funcional; APIs, consumidor e repositório materializam as variantes; carga, instrumentação e injeção de falha operacionalizam os cenários; consolidação e análise dos dados permitem discutir os trade-offs. A equivalência funcional controla a regra de processamento, mas não elimina diferenças de topologia: Kafka e consumidor adicionais integram o custo observado da variante assíncrona.

\subsection{Procedimentos técnicos}
Os procedimentos incluem estudo de artigos científicos, desenvolvimento incremental, testes unitários e de integração, calibração e experimento controlado. O levantamento bibliográfico fundamenta conceitos e situa os estudos relacionados; a resposta à pergunta de pesquisa dependerá de medições no protótipo.

\section{Procedimentos metodológicos}

\subsection{Métodos de coleta de informações}
A fundamentação utiliza exclusivamente artigos científicos publicados em periódicos e anais de conferências ou workshops. Os tipos de contribuição e as limitações dos estudos são discutidos no Capítulo~\ref{cap:revisao}.

Na coleta computacional, cada evento possui UUID e instante de ocorrência definidos antes do envio. O conjunto completo de campos é preservado em reentregas. Um identificador UUID de execução, \texttt{run\_id}, é enviado no cabeçalho \texttt{X-Run-ID} e propagado ao consumidor por cabeçalho Kafka. Esse metadado identifica a coleta, sem alterar o conteúdo usado no hash.

O gerador \texttt{scripts/load.js}, executado com k6, registra evento enviado, fase, variante, início e término da requisição, código HTTP e erro. Os logs JSONL das aplicações registram etapa, UUID, execução, resultado e marcos temporais. O marco \texttt{commit\_completed} é emitido pelo repositório comum depois do retorno da transação, tanto na variante síncrona quanto no consumidor. Uma repetição idempotente é identificada como \texttt{duplicate} e não representa uma nova linha persistida.

A coleta mantém três instantes distintos: ocorrência do fato, conclusão observada pela aplicação depois do commit e recebimento da resposta HTTP. O campo \texttt{persisted\_at} é preenchido durante a inserção no PostgreSQL; não representa o término exato do commit. A consulta \texttt{GET /audit/\{event\_id\}} verifica a visibilidade da linha, mas seu instante de resposta também não substitui o marco pós-commit.

O tempo de requisição é observado pelo gerador. Na correlação entre processos, utiliza-se a referência UTC do mesmo host; o início registrado pelo k6 tem resolução de milissegundos. A presença de timestamps em nanossegundos nos logs não aumenta a resolução da medida de ponta a ponta. Relógios monotônicos são preservados nos registros para durações locais. Antes da coleta definitiva, o piloto deve verificar resolução, consistência dos relógios e possíveis valores incompatíveis; estes serão sinalizados como falha de instrumentação, sem truncamento silencioso.

As estatísticas Docker registram CPU e memória por contêiner. Offsets do tópico e do grupo permitem observar o lag; consultas ao PostgreSQL verificam existência, hash e unicidade. Os scripts preservam arquivos brutos, configuração, versões, identificação das imagens e checksums. O estado de validação desses instrumentos é apresentado na Seção~\ref{sec:implementacao-real}.

\subsection{Métodos de análise das informações}
A análise partirá dos dados brutos preservados, correlacionados por execução e UUID. Para cada repetição serão calculados quantidade oferecida, aceita e concluída; média, mediana, desvio-padrão e percentis p95 e p99 das latências; throughput concluído; erros por categoria; backlog máximo; reentregas e recuperação. A taxa oferecida efetiva e as iterações que o gerador não conseguiu iniciar serão apresentadas separadamente.

Primeiro serão examinados os resultados de cada repetição. A agregação utilizará repetições independentes, acompanhadas de dispersão, tamanho de efeito e intervalos de confiança quando os dados sustentarem essas estimativas. A opção entre técnicas paramétricas e não paramétricas dependerá dos pressupostos e será justificada. Testes com poucos eventos destinados à depuração não sustentarão inferências sobre desempenho ou superioridade arquitetural.

\section{Metodologia de desenvolvimento do sistema}

\subsection{Modelo de desenvolvimento de software}
O desenvolvimento é incremental: contrato, serialização e persistência; variante síncrona; produtor e consumidor; tratamento de falhas; instrumentação; calibração; coleta definitiva. Cada incremento mantém requisitos, código, diagramas e verificações correspondentes. Casos com doubles de bibliotecas verificam decisões locais; testes de integração utilizam Kafka e PostgreSQL reais.

\subsection{Tecnologias utilizadas}
O protótipo utiliza Python 3.12, FastAPI, PostgreSQL 16, Apache Kafka 3.9.1 e Docker Compose. O contrato declarativo e o repositório comum reduzem divergências funcionais. PostgreSQL oferece transações e unicidade; Kafka oferece retenção e posições de consumo; Compose descreve a implantação local. k6 1.4.0 e scripts Python compõem a geração e consolidação dos dados.

Cada execução deve registrar sistema operacional, processador, memória, versões efetivas, identificadores das imagens, limites de CPU e memória, partições, replicação, parâmetros de timeout e retry e propriedades da carga. O snapshot dos arquivos é identificado por SHA-256; um hash de commit será acrescentado quando houver repositório Git. Credenciais permanecem na configuração privada do ambiente e não integram relatórios, capturas ou dados experimentais.

\subsection{Arquitetura do sistema}
As variantes compartilham \texttt{AuditEvent}, serialização determinística, SHA-256 e repositório PostgreSQL. Na síncrona, a resposta 201 sucede a transação ou o reconhecimento de duplicata. Na assíncrona, a API aguarda confirmação individual de publicação com \texttt{acks=all}; o consumidor processa posteriormente. A resposta 202 confirma essa publicação, sem atestar persistência no PostgreSQL. Um timeout de publicação é classificado como resultado incerto, pois a mensagem pode ter sido entregue sem que a confirmação chegasse ao cliente.

O consumidor desabilita commit automático de offsets. Após inserção ou reentrega idêntica, confirma a posição consumida. Erros transitórios de persistência recebem tentativas limitadas com espera exponencial; mensagens inválidas, conflitos ou tentativas esgotadas são encaminhados à DLQ. A posição de origem somente é confirmada depois do ACK da DLQ. Se a publicação na DLQ ou a confirmação da posição falhar, a sessão é retomada sem avançar além da mensagem não resolvida. Isso permite reentrega e exige idempotência; não estabelece uma transação única entre Kafka e PostgreSQL.

A topologia local contém um broker e uma réplica por tópico. A confirmação configurada não demonstra tolerância à perda física do broker nem alta disponibilidade de cluster. Limites de recursos do ambiente de desenvolvimento estão explícitos no Compose; sua adequação às configurações definitivas deverá ser verificada no piloto.

\section{Procedimentos de teste e validação}

\subsection{Tipos de testes aplicados}
Testes unitários verificam contrato, serialização, hash, espera de ACK, classificação de erros, retry, ordem entre DLQ e confirmação da origem e respostas HTTP. Testes de integração verificam publicação, consumo, persistência, conflitos e submissões concorrentes com serviços reais. Testes funcionais de falha interrompem dependências do ambiente isolado, verificam retomada e preservam os registros de execução. Seus resultados documentam o software e não respondem isoladamente à pergunta experimental.

O piloto verificará geração de carga, coleta, resolução temporal, limites e duração. Correções de instrumentação serão permitidas durante essa fase. O protocolo será congelado antes da coleta definitiva; uma alteração posterior relevante exigirá nova identificação de configuração e repetição das coletas afetadas.

\subsection{Cenários de teste}
C1 utilizará carga baixa para observar correção sem saturação. C2 aplicará níveis graduais para avaliar latência e capacidade. C3 produzirá uma rajada acima da taxa estável. C4 introduzirá indisponibilidade durante entrada ativa. C5 acompanhará a retomada e a drenagem do backlog.

O piloto estimará uma taxa de referência estável por cinco minutos, sem crescimento contínuo de backlog nem erros de saturação. As proporções iniciais de C2 são 25\%, 50\%, 75\% e 100\%; C1 utiliza o menor nível. Para preservar a equivalência de volume, cada par comparado receberá a mesma taxa absoluta e o mesmo perfil de eventos. A escolha da referência comum, da intensidade e duração da rajada e da tolerância à diferença entre taxa programada e efetiva será registrada antes da coleta.

A definição operacional de C4 permanece uma decisão prévia necessária: interromper o processo que persiste em cada variante (API síncrona e consumidor) avalia indisponibilidade desses processos; interromper o PostgreSQL em ambas avalia uma dependência comum. São falhas diferentes e não devem ser agregadas sob o mesmo identificador sem distinção. Também devem ser fixados duração da interrupção, prazo de drenagem e critério de estabilidade de C5.

Cada execução definitiva terá 60 segundos de aquecimento excluídos da análise e 300 segundos de medição, salvo mudança justificada pelo piloto e registrada antes do congelamento. Serão realizadas no mínimo dez repetições por combinação de variante e cenário. A ordem das variantes será alternada e randomizada em blocos. Entre repetições definitivas, apenas o estado do projeto experimental isolado será reinicializado, com verificação posterior de serviços, banco e tópicos. Ensaios curtos de validação preservam identificação própria e não são contados entre essas repetições.

\subsection{Métricas coletadas}
A latência HTTP corresponde ao tempo entre início do envio e recebimento da resposta. Na variante síncrona, a resposta de sucesso sucede a persistência; na assíncrona, sucede o ACK de publicação. A latência de conclusão utiliza, em ambas as variantes, o início do envio e o marco comum pós-commit. Os dois indicadores serão apresentados separadamente, evitando comparar aceitação de uma variante com conclusão da outra.

Throughput concluído é a quantidade de UUIDs únicos efetivamente persistidos na janela de medição dividida pela duração dessa janela. Eventos concluídos somente durante a drenagem são contabilizados à parte. Entradas oferecidas, respostas aceitas e linhas persistidas constituem conjuntos distintos. Uma confirmação pós-commit cujo registro de envio esteja ausente é uma inconsistência de coleta, e não uma amostra válida de latência.

Taxa de erro será discriminada em HTTP, publicação, validação, persistência, conflito e falha terminal. O lag de Kafka resulta da diferença entre posição final e posição confirmada por partição. A diferença entre eventos aceitos e persistidos constitui outro indicador e não será tratada como sinônimo de lag. Tentativas locais de retry, releituras da mesma posição Kafka e reenvios de um mesmo UUID também serão contados separadamente.

Ao fim do prazo de observação, os eventos aceitos sem persistência serão separados em falhas terminais registradas na DLQ, eventos ainda pendentes e ausências não conciliadas. Uma mensagem em DLQ foi retida para inspeção; não representa sucesso de auditoria, mas tampouco demonstra desaparecimento no transporte. A afirmação de perda exige investigar a ausência na persistência e nas evidências de retenção. Submissões de resultado incerto serão conciliadas sem serem automaticamente classificadas como aceites confirmados ou perdas.

Duplicidade persistida corresponde a mais de uma linha para o mesmo UUID e deve ser zero. Conflito corresponde ao mesmo UUID associado a conteúdo diferente. O tempo de recuperação será medido desde a retomada até o atendimento do critério de estabilidade de backlog e throughput fixado no protocolo. O tempo até lag zero observado em uma validação curta não substitui esse critério científico.

\subsection{Critérios de aceitação}
O núcleo será considerado funcional nos cenários verificados quando: eventos válidos forem persistidos; entradas inválidas forem rejeitadas; 201 ocorrer após confirmação da transação; 202 depender de ACK individual; duplicatas preservarem uma linha; conflitos preservarem o registro original; retry e DLQ obedecerem à política; e offsets de origem não forem confirmados antes da persistência ou retenção na DLQ.

A coleta definitiva será aceita se configuração e snapshot forem registrados, instrumentos estiverem validados, protocolo estiver congelado, aquecimento estiver separado, repetições mínimas forem concluídas, ordem estiver registrada e dados brutos permitirem conciliação. Falhas do instrumento, reinicializações externas e dados incompletos serão registrados como causas de invalidade conforme critérios prévios. Um resultado desfavorável à hipótese não invalida a execução.

\subsection{Protocolo de execução e reprodutibilidade}
O executor prepara a configuração privada; constrói as imagens; inicia o projeto Compose isolado; aguarda os testes de prontidão; confirma esquema e tópicos; identifica o snapshot; executa aquecimento e carga; registra os instantes de falha e retomada, quando aplicável; aguarda drenagem; exporta logs, offsets e registros; calcula checksums; e consolida indicadores a partir dos artefatos.

Cada \texttt{run\_id} possui manifesto, eventos enviados, saída do k6, logs JSONL, estatísticas de recursos, offsets e exportação do banco. Diretórios de execuções anteriores não são sobrescritos. A documentação da implementação apresenta os comandos e distingue o modo de validação do protocolo definitivo. Erros de execução permanecem registrados com seu motivo.

\subsection{Ameaças à validade}
A validade interna é afetada por concorrência com outros processos, caches, aquecimento, virtualização, limites de recursos, coleta de telemetria e ordem dos ensaios. A mitigação inclui registrar o ambiente, utilizar configurações estáveis, reduzir processos concorrentes durante a coleta, repetir e alternar a ordem. O custo do broker adicional pertence ao objeto comparado; por isso, a pesquisa não isola universalmente um estilo arquitetural.

A validade de construto depende dos marcos temporais e da classificação dos eventos. Separar ACK, commit e resposta HTTP reduz ambiguidade. Resolução de relógios, atraso da observação e custo de emissão dos logs limitam a precisão. Lag, fila pendente, DLQ e perda são estados distintos, cujo tratamento deve ser verificável nos dados.

A validade de conclusão depende de repetições suficientes e de sua independência. A análise examinará distribuições, intervalos e magnitudes, sem interpretar milhares de eventos correlacionados como milhares de experimentos. Resultados de validação funcional não serão extrapolados para desempenho.

A validade externa é limitada por dados sintéticos, domínio simplificado, uma implementação, broker único, versões e ambiente local. Não serão generalizadas conclusões para todas as APIs REST ou arquiteturas EDA. A preservação do protocolo, das versões e dos dados permite replicações que investiguem essas limitações.
